\section*{अध्याय \ref{ch:g6:circuits-and-safety} --- विद्युत परिपथ और सुरक्षा}

\begin{solution}{exo:g6:circuits-and-safety:1}
ईमानदारी से: जोड़ --- कौन-सा उपकरण किससे जुड़ा है, और सड़कें कहाँ
फूटती हैं। अनदेखा: असली तारों की लंबाई, रंग और मेज़ पर उनका \omterm{def:g2:pushes-and-pulls:force}{बल} खाता
रास्ता।
\end{solution}

\begin{solution}{exo:g6:circuits-and-safety:2}
\omterm{def:g3:first-electric-circuit:battery}{बैटरी}: अलग-अलग लंबाई की दो समांतर लकीरें; लैंप: क्रॉस वाला वृत्त;
खुला \omterm{ex:g3:first-electric-circuit:switch}{स्विच}: उठा हुआ उठाऊ पुल; बंद \omterm{ex:g3:first-electric-circuit:switch}{स्विच}: बिछा हुआ वही पुल; \omterm{ex:g6:circuits-and-safety:motor}{मोटर}:
वृत्त में M; \omterm{ex:g6:circuits-and-safety:motor}{बज़र}: चोट खाया वृत्त। \omterm{def:g3:first-electric-circuit:battery}{बैटरी} का $+$ \emph{लंबी} लकीर बताती
है।
\end{solution}

\begin{solution}{exo:g6:circuits-and-safety:3}
एक आयत: \omterm{def:g3:first-electric-circuit:battery}{बैटरी}, दबाने वाला \omterm{ex:g3:first-electric-circuit:switch}{स्विच} और \omterm{ex:g6:circuits-and-safety:motor}{बज़र}, सब एक ही फंदे में। उँगली
हटते ही \omterm{ex:g3:first-electric-circuit:switch}{स्विच} खुल जाता है --- यानी इकलौती सड़क में दरार, और फंदे का
नियम \omterm{ex:g6:circuits-and-safety:motor}{बज़र} को उसी पल चुप कर देता है।
\end{solution}

\begin{solution}{exo:g6:circuits-and-safety:4}
ज़रूरी नहीं। फ़ैसला करने वाला सवाल: कौन किससे जुड़ा है? उपकरण वही और
जोड़ वही --- तो परिपथ भी वही, उलझनें चाहे कितनी ही अलग दिखें।
\end{solution}

\begin{solution}{exo:g6:circuits-and-safety:5}
पंखा घूमता रहता है: लैंप की मौत सिर्फ़ उसकी अपनी \omterm{def:g5:series-parallel-circuits:parallel}{शाखा} अँधेरी करती है
--- समांतर वाली स्वतंत्रता। और खुला \omterm{ex:g3:first-electric-circuit:switch}{स्विच} फूट से पहले साझा सड़क पर
बैठा है, इसलिए वह दोनों को रोक देता है: हुक्म की चौकी वाला नियम।
\end{solution}

\begin{solution}{exo:g6:circuits-and-safety:6}
कोई नंगी \omterm{def:g4:conductors-insulators:conductor}{चालक} सड़क जो स्रोत के दोनों पक्षों को सीधे जोड़ दे और हर
उपकरण को छोड़ दे: बिना रुकावट धारा दौड़ पड़ती है और तार तपते हैं।
शिकार: \omterm{def:g3:first-electric-circuit:battery}{बैटरी} (चुककर बरबाद) और तार की \omterm{def:g4:conductors-insulators:insulator}{विद्युतरोधी} परत --- और घर की
बिजली में कभी-कभी पूरा घर।
\end{solution}

\begin{solution}{exo:g6:circuits-and-safety:7}
उसने पलक झपकते अपनी \omterm{def:g5:series-parallel-circuits:parallel}{शाखा} काट दी, क्योंकि धारा दौड़ने लगी थी ---
यानी कोई ख़राबी, कोई सनक नहीं। उसे दोबारा चढ़ाने से पहले ख़राबी ढूँढ़कर
ठीक करो (शक वाला उपकरण निकालो, \omterm{def:g5:series-parallel-circuits:parallel}{शाखा} देखो)।
\end{solution}

\begin{solution}{exo:g6:circuits-and-safety:8}
नहाने की जगह: नम त्वचा \omterm{def:g4:conductors-insulators:conductor}{चालक} है, और गीले शरीर के रास्ते घर की बिजली का
धक्का जानलेवा हो सकता है। ख़राब तार: चटकी हुई परत दरार वाली बाड़ है ---
उँगलियाँ ताँबे से मिल सकती हैं। \omterm{def:g3:first-electric-circuit:battery}{बैटरी} हाँ, सॉकेट कभी नहीं: \omterm{def:g3:first-electric-circuit:battery}{बैटरी} का
धक्का बहुत हल्का है; सॉकेट का सैकड़ों गुना ताक़तवर।
\end{solution}

\begin{solution}{exo:g6:circuits-and-safety:9}
आरेख: \omterm{def:g3:first-electric-circuit:battery}{बैटरी}, साझा सड़क पर बंद \omterm{ex:g3:first-electric-circuit:switch}{स्विच}, फिर एक फूट --- \omterm{ex:g6:circuits-and-safety:motor}{बज़र} वाली \omterm{def:g5:series-parallel-circuits:parallel}{शाखा} और
लैंप वाली \omterm{def:g5:series-parallel-circuits:parallel}{शाखा} --- जो मिलकर $-$ तक लौटती हैं। अंदाज़ा: लैंप जलता है
\emph{और} \omterm{ex:g6:circuits-and-safety:motor}{बज़र} बजता है, दोनों आज़ाद। \omterm{ex:g6:circuits-and-safety:motor}{बज़र} की \omterm{def:g5:series-parallel-circuits:parallel}{शाखा} पर अतिरिक्त \omterm{ex:g3:first-electric-circuit:switch}{स्विच} खुला
हो तो: वहाँ सन्नाटा, पर लैंप जलता रहता है --- \omterm{def:g5:series-parallel-circuits:parallel}{शाखा} के \omterm{ex:g3:first-electric-circuit:switch}{स्विच} सिर्फ़
अपनी \omterm{def:g5:series-parallel-circuits:parallel}{शाखा} को हुक्म देते हैं।
\end{solution}

\begin{solution}{exo:g6:circuits-and-safety:10}
अतिरिक्त तार एक \omterm{def:g6:circuits-and-safety:short}{लघुपथन} है: $+$ से $-$ तक सीधी नंगी सड़क। धारा मेहनत
कराने वाली \omterm{ex:g6:circuits-and-safety:motor}{मोटर} के बजाय आसान ख़ाली सड़क पसंद करती है, इसलिए पंखा भूखा
रह जाता है जबकि वह छोटा रास्ता पूरी दौड़ ढोता है और तपता है ---
\omterm{def:g3:first-electric-circuit:battery}{बैटरी} चुक जाती है और जलने का न्योता मिल जाता है। ``मज़बूती'' वाला तार
हटा दो।
\end{solution}

\begin{solution}{exo:g6:circuits-and-safety:11}
दीवार का \omterm{ex:g3:first-electric-circuit:switch}{स्विच} फ़िटिंग के अपने फंदे को खोल देता है --- लैंप की सड़क
कट जाती है। विच्छेदक और पीछे जाकर पूरी \omterm{def:g5:series-parallel-circuits:parallel}{शाखा} खोल देता है --- यानी
फ़िटिंग की तारबंदी ख़राब हो या दीवार के \omterm{ex:g3:first-electric-circuit:switch}{स्विच} पर ग़लत लेबल लगा हो, तब
भी \omterm{def:g5:series-parallel-circuits:parallel}{शाखा} मुर्दा रहती है। \omterm{def:g4:series-circuits:series}{श्रेणी में} लगी दो दरारें: कोई एक भी धारा रोक
देती है, और दोनों मिलकर किसी एक के बारे में हुई ग़लती माफ़ कर देती
हैं।
\end{solution}

\begin{solution}{exo:g6:circuits-and-safety:12}
\omterm{def:g3:first-electric-circuit:battery}{बैटरी}, फिर साझा सड़क पर मुख्य \omterm{ex:g3:first-electric-circuit:switch}{स्विच}, और फिर चार समांतर शाखाएँ: मंच का
लैंप 1; मंच का लैंप 2; अपने श्रेणी वाले \omterm{ex:g3:first-electric-circuit:switch}{स्विच} के साथ परदे की \omterm{ex:g6:circuits-and-safety:motor}{मोटर}; और
अपने श्रेणी वाले \omterm{ex:g3:first-electric-circuit:switch}{स्विच} के साथ \omterm{ex:g6:circuits-and-safety:motor}{बज़र}। जाँच: मुख्य \omterm{ex:g3:first-electric-circuit:switch}{स्विच} सब पर राज करता
है (साझा सड़क); लैंप एक-दूसरे से आज़ाद रहकर ख़राब हो सकते हैं (अलग
शाखाएँ); \omterm{ex:g6:circuits-and-safety:motor}{मोटर} और \omterm{ex:g6:circuits-and-safety:motor}{बज़र} अपने-अपने \omterm{ex:g3:first-electric-circuit:switch}{स्विच} की मानते हैं (अपनी \omterm{def:g5:series-parallel-circuits:parallel}{शाखा} के भीतर
श्रेणी) और एक-दूसरे को अनदेखा करते हैं (\omterm{def:g5:series-parallel-circuits:parallel}{शाखाओं} के बीच समांतर)।
\end{solution}

\begin{solution}{pb:g6:circuits-and-safety:1}
\textbf{1.} \omterm{def:g4:series-circuits:series}{श्रेणी में}, एक भी मरा हुआ या बुझाया गया उपकरण सब कुछ
अँधेरा कर देता, और तीनों उपकरण धक्का बाँट लेते --- सब मंद और कमज़ोर।
समांतर हर उपकरण को आज़ाद और पूरी ताक़त पर रखता है।
\textbf{2.} हर \omterm{ex:g3:first-electric-circuit:switch}{स्विच} ठीक अपनी \omterm{def:g5:series-parallel-circuits:parallel}{शाखा} के उपकरण को हुक्म देता है; बाक़ी
\omterm{def:g5:series-parallel-circuits:parallel}{शाखाओं} पर इनमें से किसी का हुक्म नहीं चलता।
\textbf{3.} साझा सड़क पर, \omterm{def:g3:first-electric-circuit:battery}{बैटरी} और पहली फूट के बीच --- यानी हर \omterm{def:g5:series-parallel-circuits:parallel}{शाखा}
से पहले।
\textbf{4.} \omterm{def:g3:first-electric-circuit:battery}{बैटरी} $\to$ मुख्य \omterm{ex:g3:first-electric-circuit:switch}{स्विच} $\to$ \omterm{def:g5:series-parallel-circuits:parallel}{संधि}, जहाँ से तीन शाखाएँ
फूटती हैं --- (छत का लैंप $+$ उसका \omterm{ex:g3:first-electric-circuit:switch}{स्विच}), (मेज़ का लैंप $+$ उसका
\omterm{ex:g3:first-electric-circuit:switch}{स्विच}), (पंखे की \omterm{ex:g6:circuits-and-safety:motor}{मोटर} $+$ उसका \omterm{ex:g3:first-electric-circuit:switch}{स्विच}) --- जो फिर मिलकर \omterm{def:g3:first-electric-circuit:battery}{बैटरी} तक लौट
जाती हैं।
\textbf{5.} वे एक ही \omterm{def:g5:series-parallel-circuits:parallel}{शाखा} पर \omterm{def:g4:series-circuits:series}{श्रेणी में} हैं: धक्का बँट जाता है (दोनों
कमज़ोर --- छत के लैंप का मंद होना असल में मेज़ वाली जोड़ी का लक्षण
था), और पंखे का \omterm{ex:g3:first-electric-circuit:switch}{स्विच} उनकी साझा सड़क पर बैठकर दोनों को हुक्म देता है।
\textbf{6.} उन्हें अलग करो: एक नया तार डालकर मेज़ के लैंप को सीधे
आपूर्ति के दोनों सिरों पर अपनी \omterm{def:g5:series-parallel-circuits:parallel}{शाखा} दे दो --- तब हर एक पूरी ताक़त पर
और आज़ाद हो जाएगा।
\textbf{7.} सिरों के बीच \omterm{def:g6:circuits-and-safety:short}{लघुपथन}। जाँचो कि \omterm{def:g3:first-electric-circuit:battery}{बैटरी} अब भी काम देती है या
नहीं --- \omterm{def:g6:circuits-and-safety:short}{लघुपथन} उसे चुका देते हैं और सेल बरबाद कर सकते हैं (गरम हो
चुकी \omterm{def:g3:first-electric-circuit:battery}{बैटरी} सेवामुक्त कर देनी चाहिए)।
\textbf{8.} \omterm{def:g6:circuits-and-safety:short}{लघुपथन} (नंगे तारों का आपस में मिल जाना) और झटके या जलन
(उँगलियों का नंगे ताँबे से मिलना)।
\textbf{9.} नम त्वचा सूखी से कहीं बेहतर चालन करती है; \omterm{def:g3:first-electric-circuit:battery}{बैटरी} वाले जोड़
में भी सूखे हाथों का अनुशासन बनता है --- और घर में वही स्पर्श मुख्य
बिजली के सैकड़ों गुना ताक़तवर धक्के से मिलता है: शायद जानलेवा। आदतें
सुरक्षित परिपथ पर ही सधती हैं।
\textbf{10.} विच्छेदक ने पकड़ा कि उसकी \omterm{def:g5:series-parallel-circuits:parallel}{शाखा} में बहुत ज़्यादा धारा दौड़
रही थी --- कोई ख़राबी, शायद किसी फ़िटिंग में बारिश। समझदार क्रम:
विच्छेदक बंद ही रहने दो, \omterm{def:g5:series-parallel-circuits:parallel}{शाखा} के उपकरण निकालकर जाँचो, ख़राबी ठीक करो
या सुखाओ, और तभी उसे दोबारा चढ़ाओ।
\textbf{11.} अपना अलग फंदा, कोठरी के जाल से अलग (या कोठरी की \omterm{def:g3:first-electric-circuit:battery}{बैटरी} पर
एक और समांतर \omterm{def:g5:series-parallel-circuits:parallel}{शाखा} के रूप में): \omterm{def:g3:first-electric-circuit:battery}{बैटरी}, घर पर दबाने वाला \omterm{ex:g3:first-electric-circuit:switch}{स्विच}, लंबा
दुहरा तार, कोठरी पर \omterm{ex:g6:circuits-and-safety:motor}{बज़र}, और वापस --- दरवाज़े की घंटी अपने आप में एक
श्रेणी फंदा है, चाहे उसके तार कितने ही लंबे क्यों न हों।
\textbf{12.} जैसे: ``सॉकेट: हमारे प्रयोग \omterm{def:g3:first-electric-circuit:battery}{बैटरी} पर ही ख़त्म होते हैं
--- सॉकेट हमारा नहीं है। पानी: गीले हाथ कोई तारबंदी नहीं छूते, न भीतर
न बाहर। नंगे तार: हर जोड़ लपेटा हुआ; नंगा तार शॉर्टकट नहीं, ख़राबी
है।''
\end{solution}
